% Blank content skeleton based on the original thesis project.
% Fill in the empty fields and replace the chapter headings as needed.
\def\ThesisTitle{基于xxx的系统设计}
\def\CoverThesisTitle{基于xxx的系统设计}
\def\CoverHeading{本科毕业论文（设计）}
\def\TitleLabel{题　　目：}
\def\MajorLabel{年级专业：}
\def\NameLabel{姓　　名：}
\def\StudentIDLabel{学　　号：}
\def\AdvisorLabel{指导教师：}
\def\ThesisAuthor{张三}
\def\ThesisMajor{xx级电子信息工程专业}
\def\ThesisStudentID{}
\def\ThesisAdvisor{}
\def\ThesisDate{}
\def\HeaderLeft{xx届电子信息工程专业毕业论文（设计）}
\def\HeaderRight{张三：基于xxx的系统设计}
\def\ThesisAffiliation{三明学院xx级电子信息工程专业\hspace{0.5em}福建三明\hspace{0.5em}365004}
\def\ContentsTitle{目录}
\def\AbstractHeading{摘要：}
\def\KeywordsHeading{关键词：}
\def\AbstractText{针对【研究对象】在【应用场景】中出现的【关键问题】，本文围绕系统设计与性能改善开展研究。首先，分析研究对象的工作原理、应用需求及现有方法的不足，明确设计目标、功能要求与主要约束条件；其次，结合系统特点提出【核心方法】，完成总体架构设计，并对关键模块、参数选取及软硬件实现流程进行说明；随后，在【实验或仿真条件】下搭建验证平台，设计对比实验，并采用【评价指标】分析系统的稳态性能和动态响应。结果表明，【填写主要结果及结论】。研究结果可为【应用场景】中相关系统的设计与优化提供参考。}
\def\KeywordsText{研究对象；核心方法；系统设计；实验验证}
\def\ReferenceHeading{参考文献}
\def\EnglishContentsTitle{英文题目、摘要和关键词}
\def\EnglishTitle{System Design Based on xxx}
\def\EnglishAuthor{ZHANG San}
\def\EnglishAffiliation{20XX Electronic Information Engineering, Sanming University, Fujian Sanming, 365004}
\def\EnglishAbstractHeading{Abstract:\hspace{0.5em}}
\def\EnglishAbstractText{This study addresses [key problem] affecting [research subject] in [application scenario]. First, the design objectives and constraints are established based on the problem context and related work. A [core method] is then proposed, and the system architecture, key modules, and implementation process are described. Finally, the design is evaluated under [experimental or simulation conditions] using [evaluation metrics]. The results show that [insert the principal findings and conclusions].}
\def\EnglishKeywordsHeading{Keywords:\hspace{0.5em}}
\def\EnglishKeywordsText{research subject; core method; system design; experimental validation}
\def\AcknowledgmentsHeading{致谢}
\def\AcknowledgmentsText{谨向在本课题研究和论文撰写过程中给予帮助与支持的老师、同学和家人致以诚挚谢意。首先，感谢指导教师【姓名】在选题、研究思路和论文修改方面给予的指导；感谢【实验室或团队名称】的老师和同学在资料整理、实验开展及交流讨论中提供的帮助；感谢家人在学习期间给予的理解与支持。谨以此文向所有关心和帮助过我的人表示感谢。}
\def\CommitmentHeading{承诺书与授权书}
\def\CommitmentTitle{三明学院毕业论文（设计）承诺书}
\def\AuthorizationTitle{毕业论文（设计）版权使用授权书}
\def\AppendixOneHeading{附录 1}
\def\AppendixTwoHeading{附录 2}
\def\AppendixThreeHeading{附录 3}
\def\ReferenceBibFile{references/refs}
\def\ReferenceStyle{gbt7714-numerical}
\def\ReferenceNocite{}
\IncludeAppendicestrue

\long\def\BodyContent{
% The following paragraphs are drafting placeholders. Replace them with original work.
\chapter{引言}
\section{研究背景及意义}
本节介绍课题所处的应用背景，说明研究对象在实际系统中的作用，并从理论研究或工程应用角度阐述开展本课题的意义。建议结合具体问题说明研究动机，明确本文关注的研究范围。

\section{国内外研究现状}
本节围绕课题主题梳理相关研究进展。可按照研究方法、技术路线或关键问题组织文献，概括已有工作的主要思路与适用条件，并在比较分析的基础上指出仍待解决的问题。以下段落和题录仅用于演示引用及著录格式，正式稿应替换为与课题相关并已核实的文献。以变换器控制为例，已有研究涵盖全桥 DC-DC 预测控制、系统辨识及缩短预测时域的控制设计等方向\cite{Xie2009ImplicitMPC,Xie2012FullBridgeMPC,Tomesc2013ARXID,Cheng2018ReducedHorizonMPC}。后续研究还讨论了平滑模式切换、双向变换器模式过渡和参数不确定性下的自适应预测控制\cite{Kwon2018SmoothModeSwitching,Sharma2020ModeTransition,Albira2021AdaptiveMPC}。近年的研究进一步涉及自适应神经网络控制、基于 ARX 与递推最小二乘的预测控制及快速负载瞬态响应\cite{Gangula2023AdaptiveNN,Kermadi2024ARXRLS,Liu2024DualModeController,Maruta2024NNBias,Guo2025AdaptiveDuty,Xu2025BumplessTransfer,Cheng2026IncrementalDDMPC}。正式撰写时应结合文献内容进行归纳比较，不要直接保留此示例文字。
\subsection{国内研究进展}
本小节归纳国内相关研究的代表性方法与发展趋势，说明各类方法解决的问题及其适用场景，并总结与本文课题直接相关的研究基础。
\subsection{国外研究进展}
本小节归纳国外相关研究的主要方向，比较不同技术路线的特点，并结合研究对象分析其可借鉴之处和应用限制。
\section{本文主要研究内容}
本节概述本文拟解决的问题、采用的研究方法和主要工作安排。可按章节说明各部分之间的逻辑关系，使研究目标、技术路线与后续验证内容相互对应。

\chapter{理论基础与系统建模}
\section{理论基础}
本节介绍后续分析所需的基本概念、理论方法和建模假设。首次出现的专业术语、变量和缩写应给出清晰定义，并保持全文符号与单位的一致。
\section{系统模型}
本节说明研究对象的组成、边界条件及输入输出关系，逐步给出系统模型的建立过程。模型推导应列明必要假设，并说明参数来源及其物理含义。
\section{本章小结}
本章围绕研究所需的理论与模型进行了说明，为后续方案设计和结果分析提供基础。实际撰写时，可在此归纳本章得到的主要结论及其与下一章的衔接。

\chapter{系统设计与实现}
\section{总体设计}
本节从整体层面说明系统目标、功能组成及各模块之间的关系。建议配合总体框图描述信号流、能量流或数据流，并解释关键设计选择与约束条件。
\section{关键方法}
本节重点说明本文采用的关键方法或控制策略。可按照处理流程介绍方法的输入、计算步骤、输出及参数设置，并说明该方法如何对应前文提出的问题。
\section{实现过程}
本节介绍方案在硬件、软件或仿真环境中的实现过程，包括主要器件或工具、程序结构、接口关系和调试步骤。涉及具体配置时，应记录必要参数以便复现。
\section{本章小结}
本章说明了系统结构、关键方法和实现流程。正式稿中可在此归纳实现结果，并指出后续实验需要验证的功能与性能指标。

\chapter{实验结果与分析}
\section{实验平台与评价指标}
本节描述实验平台、测试设备、实验条件和数据采集方式，并定义用于评价系统性能的指标。实验条件和指标的定义应清楚、完整，保证不同方案之间可以公平比较。
\section{实验结果}
本节按照实验目的呈现测量结果，可使用图表展示典型波形、统计数据或不同方案之间的比较。图表应有编号、标题和必要说明，文中应指出结果所反映的主要现象。
\section{结果讨论}
本节结合理论分析解释实验现象，讨论结果与预期是否一致、关键因素如何影响性能，以及方法的适用范围和局限。所有结论均应由实际数据或分析过程支持。
\section{本章小结}
本章归纳实验设置、主要观察结果及其分析结论。正式稿中应突出与研究目标相关的发现，并为全文总结提供依据。

\chapter{总结与展望}
\section{工作总结}
本节回顾本文的研究问题、采用的方法和完成的主要工作，并依据前文分析与实验结果总结研究结论。总结应与论文目标相对应，避免引入正文未讨论的新结果。
\section{后续研究方向}
本节结合本文工作的适用范围和局限，提出可继续研究的问题或改进方向。展望内容应具体、可执行，并与本文结论保持关联。
}

\long\def\CommitmentContent{
我仔细阅读了三明学院毕业论文（设计）的有关文件规定。\par
我知道，抄袭别人的成果是剽窃行为，是可耻的，也是违反毕业论文（设计）规定的。如果引用别人的成果或其他公开的资料（包括网上查到的资料），必须按照规定的参考文献的表\par
述方式在正文引用处和参考文献中明确列出。\par
我郑重承诺，严格遵守学院毕业论文（设计）规定，以自己的真实水平认真做好毕业论文（设计）。如有违反规定的行为，我将接受严肃处理。我的毕业论文（设计）题目为：\par
\vspace{1.1em}
\begin{center}
学生签名：　　　　　　　　　日期：　　　年　 月　 日\\[0.5em]
\end{center}
\vspace{1.1em}
\begin{center}{\heiti\bfseries\fontsize{20pt}{28pt}\selectfont 毕业论文（设计）版权使用授权书}\par\end{center}
\vspace{0.8em}
本论文（设计）作者完全了解学校有关保留、使用论文的规定，同意学校保留并向国家有关部门或机构送交论文的复印件和电子版，允许论文（设计）被查阅和借阅。本人授权三明学\par
院可以将本论文（设计）的全部或部分内容编入有关数据库进行检索，可以采用影印、缩印或扫描等复制手段保存和汇编本论文（设计）。\par
本论文（设计）属于：保密\_\_CHECKBOX\_\_，在　　年解密后适用本授权书。\par
不保密\_\_CHECKBOX\_\_。（请在以上相应方框内打“√”）\par
\vspace{1.1em}
\begin{center}
作者签名：　　　　　　　　　日期：　　　年　 月　 日\\[0.5em]
\end{center}
}
\long\def\AppendixOneContent{请在此处填写附录内容，例如程序清单、补充推导或技术说明；无相关内容时可删除本附录。}
\long\def\AppendixTwoContent{请在此处填写附录内容，例如参数表、补充数据或计算结果；无相关内容时可删除本附录。}
\long\def\AppendixThreeContent{请在此处填写附录内容，例如图纸、实验材料或其他补充资料；无相关内容时可删除本附录。}
